% Expanded paper insert, 2026-11-03. No standalone document or bibliography.
% Source-text revision only; no PDF generation or simulation is required.
\section{Stability, Safety and Certified Comparison: Full Proofs}
\label{app:theory}
\noindent\textbf{Correspondence: Section~\ref{sec:btd_reference}.} This appendix preserves the extended model, intermediate steps and supporting interpretation behind the corresponding main-text section.


\label{detail:sec:btd_reference}

We connect the nonlinear controller of Section~\ref{detail:sec:fpi_controller}
to four separate properties: invariance of its analyzed input region,
convergence of frequency and buffer disagreements, strict buffer-capacity
safety, and positive physical clock frequency. We then establish a
strict headroom comparison with its calibrated PI baseline for one
specified initialization. The proof concerns an explicitly defined
continuous fluid reference. Callisto retains asynchronous local
sampling, delivery delays and integer frame events, which are not
silently removed from its experimental model.

\subsection{Reference assumptions and continuous-time model}
\label{detail:subsec:btd_ref_model}

\subsubsection{What the reference retains and idealizes}
The reference retains the controller of Section~\ref{detail:sec:fpi_controller}, continuous clock phases, local occupancy
feedback, the implemented PI integrate-before-output convention and the proposed
incremental normalized-min Fuzzy surface. It assumes zero physical link latency,
zero control-execution delay, fluid rather than integer occupancy and a
common physical sampling grid $t_k=kH$, $H>0$. A zero-order hold (ZOH)
means that the applied frequency correction is constant on each interval
$[t_k,t_{k+1})$. These are explicit analysis assumptions; local-clock
polling in Callisto is not asserted to have this common grid.
The equations sampled at $kH$ below are exact endpoint relations for that
continuous ZOH flow. They are neither numerical Euler integration nor a
discrete firing-count replacement of the plant.

Let $\mathcal N_i^-$ be the fixed incoming-neighbor set and
$d_i=|\mathcal N_i^-|$. The graph is connected, undirected and unit-weight,
with a receiving channel in both directions. Its symmetric Laplacian
$L$ has diagonal $d_i$ and off-diagonal $-1$ at neighboring nodes. Each
undirected connection contributes once to this Laplacian; the pair of
directed receiving channels does not double its weights.

The physical phase is $\theta_i$, the uncontrolled frequency is
$\omega_i^u$, the applied correction is $u_i$, and the requested correction
is $q_i$. The target occupancy of receiving edge $j\to i$ is
$\beta^*_{j\to i}$. The fixed coefficient scale in the PI baseline law is
$h=p/\omega_{\rm base}$, where $p$ is a local polling period and
$\omega_{\rm base}$ is the implementation reference frequency; $H$ is the
physical sampling period imposed only in this reference. The integral
coefficient is $\kappa_I=k_i h$. These two time scales are kept distinct.
All receiving buffers have common capacity $C>0$, and the uncontrolled
frequencies $\omega_i^u$ are constant. The all-ones column vector is $\mathbf1$.

The transpose is denoted by $\top$. A symmetric matrix $A\succeq0$ is
positive semidefinite; $A\succ0$ is positive definite. The matrix order
$A\preceq B$ means $B-A\succeq0$. The symbol $\otimes$ denotes the
Kronecker product, and $\|z\|_\infty=\max_i|z_i|$.

\subsubsection{Compatible occupancy targets and marking}
Throughout the proved centered reference, all initial phases are equal.
For zero delay the fluid law in the additive phase-offset convention is
\begin{equation}
 \beta_{j\to i}(t)=\theta_j(t)-\theta_i(t)+\lambda_{j\to i}.
 \label{detail:eq:btd_d_fluid}
\end{equation}
Here $\lambda_{j\to i}$ is the fixed additive phase offset. With equal
initial phases it also equals the initial marking; this presentation does
not identify these conventions for arbitrary unequal initial phases.
The law follows the bittide fluid modeling framework
\cite{lall2024logicalSynchrony,lall2022resistance}. Define
$\beta^*_{j\to i}=\theta_j(0)-\theta_i(0)+\lambda_{j\to i}$, so the
reference starts at its compatible targets. For a common constant ramp
frequency $\omega_{\rm ref}$ put
\begin{equation}
 x_k=\theta(t_k)-\theta(0)-t_k\omega_{\rm ref}\mathbf1.
 \label{detail:eq:btd_ref_state}
\end{equation}
The common ramp cancels from edge differences. Subtracting the target in
\eqref{detail:eq:btd_d_fluid} gives
\begin{equation}
 \beta_{j\to i}(t_k)-\beta^*_{j\to i}=x_{j,k}-x_{i,k},\qquad E_k=-Lx_k.
 \label{detail:eq:btd_buffer_phase}
\end{equation}
The vector $E$ collects the Section~\ref{detail:sec:fpi_controller} frame errors, actual minus target.
Not every set of per-edge targets is compatible with a fixed marking and
phase initialization; that compatibility is an assumption of this model.

For the worked example, equal initial phases and target 1024 imply
$\lambda_{j\to i}=1024$. The delayed Callisto simulations have marking 1044 with
20 initial in-flight frames; simply retaining 1044 while imposing zero
delay and the same centered occupancy would change the initialization.

\subsection{From the original controller to the sampled recurrence}

Define the uncontrolled frequency offsets $\delta\omega=\omega^u-\omega_{\rm ref}\mathbf1$
and the physical frequency offsets $w_k=\delta\omega+u_k$. For $k\ge1$,
let $v_k=Hw_{k-1}$. The plant's previous genuine interval gives
\begin{equation}
 x_k-x_{k-1}=Hw_{k-1}=v_k.
 \label{detail:eq:btd_d_plant}
\end{equation}
Indeed, on each zero-order-held interval,
\begin{equation}
 x(t)=x_k+(t-t_k)w_k,\qquad t\in[t_k,t_{k+1}],
 \label{detail:eq:btd_exact_continuous_flow}
\end{equation}
where $x(t)=\theta(t)-\theta(0)-t\omega_{\rm ref}\mathbf1$.
Integration gives $x_{k+1}=x_k+Hw_k$ exactly. The controller is updated
only at the endpoints, but the plant evolves continuously throughout
the interval.

The factor $H$ is physical elapsed time; the controller factor $h$ remains
the fixed PI coefficient scale. Section~\ref{detail:sec:fpi_controller} uses
$e_i=-\gamma_i E_i$, so the Fuzzy inputs are
\begin{equation}
 X_k=K_e\Gamma Lx_k,\qquad
 Y_k=K_{\Delta e}\Gamma Lv_k,\qquad
 \Gamma=\operatorname{diag}(\gamma_1,\ldots,\gamma_n).
 \label{detail:eq:btd_d_inputs}
\end{equation}
Here $K_e,K_{\Delta e}$ are chosen common for readability; node-dependent
output scaling can still produce common calibrated $k_p,k_i$.

On $\Omega_r=\{(X,Y):|X|+|Y|\le r\}$, $0<r\le5$, Section~\ref{detail:sec:fpi_controller} gives
$F=\alpha F_0$, where $F_0=(X+Y)/2$ and
$\underline\alpha=(1+r/5)^{-1}\le\alpha\le\overline\alpha=2$.
With $D_{\alpha,k}=\operatorname{diag}(\alpha_{1,k},\ldots,\alpha_{n,k})$,
the calibrated increment becomes
\begin{equation}
 u_k-u_{k-1}=-D_{\alpha,k}(\kappa_I Lx_k+k_pLv_k).
 \label{detail:eq:btd_d_command_increment}
\end{equation}
Multiplication by $H$ and the next plant interval yield
\begin{align}
 v_{k+1}&=v_k-D_{\alpha,k}(bLx_k+aLv_k),\nonumber\\
 x_{k+1}&=x_k+v_{k+1},\qquad
 a=Hk_p,\quad b=H\kappa_I=Hk_i h.
 \label{detail:eq:btd_full_recurrence}
\end{align}
The extra $H$ in $b$ is essential: $v$ is a scaled frequency. The surface
factor $\alpha_{i,k}$ is node-dependent and state-dependent; freezing it
and checking one constant-gain PI characteristic is insufficient.

\subsubsection{The initial constructor call}
The controller's stored previous error is not a previous plant interval.
For the centered initialization $x_0=0$ with zero PI sample sum, zero
Fuzzy accumulator and zero previous percentage error, the first original
controller call gives $u_0=0$. Advancing $[0,H]$ then gives
\begin{equation}
 x_1=H\delta\omega,\qquad v_1=H\delta\omega.
 \label{detail:eq:btd_bootstrap}
\end{equation}
Only from $k=1$ does~\eqref{detail:eq:btd_d_plant} justify the recurrence.
Nonzero initial occupancy error is outside this centered formulation:
one must also introduce target-compatible phase coordinates or retain the
affine error bias, then explicitly evaluate the original constructor and
recheck the initial state and certificate. Evaluating the first call alone
does not restore $E=-Lx$. The centered bootstrap is consistent with, but separate from, the
event-level first-call PI/$F_0$ equivalence proved in Section~\ref{detail:sec:fpi_controller}.

\subsection{Disagreement projection and the symmetric sector matrix}

Since $L\mathbf1=0$, a common phase shift does not change occupancies.
Choose $U$ with orthonormal columns spanning the subspace orthogonal to
$\mathbf1$, so $U^\top U=I_{n-1}$ and $U^\top\mathbf1=0$. Choose its
columns as graph eigenvectors, giving
\begin{equation}
 L=U\Lambda U^\top,\qquad
 \Lambda=\operatorname{diag}(\mu_2,\ldots,\mu_n),\quad 0<\mu_2\le\mu_n.
 \label{detail:eq:btd_d_spectrum}
\end{equation}
$I_{n-1}$ denotes the identity matrix of dimension $n-1$. Define
\begin{equation}
 \xi_k=\Lambda^{1/2}U^\top x_k,\qquad
 \eta_k=\Lambda^{1/2}U^\top v_k,\qquad y_k=(\xi_k^\top,\eta_k^\top)^\top.
 \label{detail:eq:btd_projection}
\end{equation}
The symbol $\xi$ is projected phase and $\eta$ already contains the
frequency scale $H$. The coordinates retain the physical-time scale $H$; in particular,
$\eta$ is a scaled frequency disagreement rather than an integral-controller state.

For $M_k=\Lambda^{1/2}U^\top D_{\alpha,k}U\Lambda^{1/2}$,
substitution in~\eqref{detail:eq:btd_full_recurrence} gives
\begin{align}
 \eta_{k+1}&=-bM_k\xi_k+(I_{n-1}-aM_k)\eta_k,\nonumber\\
 \xi_{k+1}&=(I_{n-1}-bM_k)\xi_k+(I_{n-1}-aM_k)\eta_k.
 \label{detail:eq:btd_d_projected_recurrence}
\end{align}
In particular, the diagonal $D_\alpha$ need not commute with $L$; the
projection produces the symmetric $M_k$ instead of assuming shared modes.
Its quadratic form satisfies
\begin{equation}
 \underline\alpha\Lambda\preceq M_k\preceq\overline\alpha\Lambda,
 \qquad \operatorname{eig}(M_k)\subseteq
 [\nu_-,\nu_+]=[\underline\alpha\mu_2,\overline\alpha\mu_n].
 \label{detail:eq:btd_d_sector_matrix}
\end{equation}
The scalar $\nu$ used below is a possible eigenvalue of $M$, not a node
frequency or the number of communication edges.

\subsection{A common Lyapunov certificate from two endpoints}
\label{detail:subsec:btd_ref_certificate}

For matrices $M\in\mathbb R^{(n-1)\times(n-1)}$ and scalars $\nu$,
define
\begin{equation}
 \mathcal B(M)=\begin{bmatrix}I-bM&I-aM\\-bM&I-aM\end{bmatrix}.
 \label{detail:eq:btd_matrix_dynamics}
\end{equation}
The identity blocks have dimension $n-1$. Equations above give
$y_{k+1}=\mathcal B(M_k)y_k$. A scalar block associated with an eigenvalue
of $M_k$ is
\begin{equation}
 B(\nu)=\begin{bmatrix}1-b\nu&1-a\nu\\-b\nu&1-a\nu\end{bmatrix}.
 \label{detail:eq:btd_scalar_block}
\end{equation}
For a symmetric positive-definite matrix $\bar P\in\mathbb R^{2\times2}$,
put $P=\bar P\otimes I_{n-1}$. The scalar
$V_k=y_k^\top P y_k$ is a quadratic measure of phase and frequency
disagreement, not of absolute phase.

\begin{btdlemma}[Endpoint certificate]
\label{detail:lem:btd_endpoint}
Suppose $0<\rho<1$ and
\begin{align}
 B(\nu_-)^\top\bar P B(\nu_-)&\preceq\rho\bar P,\nonumber\\
 B(\nu_+)^\top\bar P B(\nu_+)&\preceq\rho\bar P.
 \label{detail:eq:btd_d_endpoints}
\end{align}
Then every update whose inputs satisfy the regional sector obeys
\begin{equation}
 V_{k+1}\le\rho V_k.
 \label{detail:eq:btd_contraction}
\end{equation}
The conclusion permits node-dependent and time-varying sector factors.
\end{btdlemma}
\begin{proof}
First consider one possible scalar eigenvalue. For
$\nu=(1-t)\nu_-+t\nu_+$, $0\le t\le1$, affine dependence gives
$B(\nu)=(1-t)B_-+tB_+$, where $B_\pm=B(\nu_\pm)$.
Expanding the two quadratic expressions yields
\begin{align}
 &(1-t)B_-^\top\bar P B_-+tB_+^\top\bar P B_+
       -B(\nu)^\top\bar P B(\nu)\nonumber\\
 &\qquad=t(1-t)(B_+-B_-)^\top\bar P(B_+-B_-)\succeq0.
 \label{detail:eq:btd_d_convexity}
\end{align}
Together with the endpoint assumptions, this proves
$B(\nu)^\top\bar P B(\nu)\preceq\rho\bar P$ throughout the
continuous interval $[\nu_-,\nu_+]$. No eigenvalue grid is required.

Next, each $M_k$ is real symmetric, so it admits an orthogonal
factorization $M_k=Q_k\operatorname{diag}(\nu_{1,k},\ldots,
\nu_{n-1,k})Q_k^\top$. Rotating both $\xi_k$ and $\eta_k$ by $Q_k^\top$
reduces the dynamics to the scalar blocks $B(\nu_{\ell,k})$.
The corresponding block rotation leaves $P=\bar P\otimes I_{n-1}$
unchanged. Thus summing the scalar inequalities gives
$\mathcal B(M_k)^\top P\mathcal B(M_k)\preceq\rho P$, which is
\eqref{detail:eq:btd_contraction}. Because the witness $P$ is the same for
every permitted $M_k$, $Q_k$ may change at every step. No simultaneous
diagonalization of the node-dependent matrices is assumed.
\end{proof}

The remaining task is to show that the actual nonlinear inputs stay in
the region where the sector was obtained. Stability of a family of
hypothetical linear blocks alone would not establish that fact.

\subsection{Closing the invariant input region}

Let $\mathbf e_i$ denote the $i$th unit vector and define the row
$a_i^\top=\mathbf e_i^\top U\Lambda^{1/2}$. Its squared Euclidean norm
is $\|a_i\|^2=L_{ii}=d_i$. The Fuzzy inputs become
\begin{equation}
 X_{i,k}=\gamma_iK_ea_i^\top\xi_k,\qquad
 Y_{i,k}=\gamma_iK_{\Delta e}a_i^\top\eta_k.
 \label{detail:eq:btd_d_linear_outputs}
\end{equation}
For any vector $c$ and $P\succ0$, Cauchy--Schwarz in the quadratic
metric gives $(c^\top y)^2\le(y^\top P y)(c^\top P^{-1}c)$.
Also $|X|+|Y|=\max_{\sigma=\pm1}|X+\sigma Y|$. Applying that output
bound on $V\le R$ yields
\begin{align}
 |X_{i,k}|+|Y_{i,k}|&\le\sqrt{R\Psi_i},\nonumber\\
 \Psi_i&=\gamma_i^2d_i\max_{\sigma=\pm1}
 [K_e,\sigma K_{\Delta e}]\bar P^{-1}
 [K_e,\sigma K_{\Delta e}]^\top.
 \label{detail:eq:btd_input_factor}
\end{align}
The scalar $\Psi_i$ is the largest squared input gain on the energy
ellipsoid. If $V_1\le R$ and $R\Psi_i\le r^2$ for every node, then:
each node's first scaled input pair belongs to $\Omega_r$; the valid sector gives
$V_2\le\rho V_1\le R$; the same output bound keeps the next inputs in
$\Omega_r$; induction repeats. This proves the sector domain and
disagreement ellipsoid invariant together, rather than assuming either.

Equivalently, a valid radius must satisfy
\begin{equation}
 V_1\le R\le\min_i\frac{r^2}{\Psi_i}.
 \label{detail:eq:btd_sector_radius_condition}
\end{equation}
The initialization, normalization through $d_i$ and $C$, the Fuzzy input
scales and the graph-dependent witness therefore enter the same
condition. When $\Psi_i=0$, that node places no restriction on $R$;
otherwise the displayed quotient is finite. The inference surface is
still nonlinear inside the invariant region. The sector is an enclosure
of its instantaneous response, not an assertion that its gain is fixed.

Because $P\succ0$ and $V_k\le\rho^{k-1}V_1$, the projected phase and
frequency decay exponentially. The ellipsoid constrains disagreement;
the common physical phase continues to grow and is not bounded by $V$.
Thus convergence means agreement of rates and restoration of compatible
occupancy targets, not convergence of the absolute clock phases to constants.

\subsection{From disagreement to physical buffer and clock bounds}

\subsubsection{Edge occupancy and intersample time}
For edge $e=(j\to i)$, define its signed unit incidence vector
$\mathbf b_e=\mathbf e_j-\mathbf e_i$. Since it is orthogonal to
$\mathbf1$, its output from the projected phase is
\[
 \mathbf b_e^\top x_k=\mathbf b_e^\top U\Lambda^{-1/2}\xi_k.
\]
Define the graph pseudoinverse $L^\dagger=U\Lambda^{-1}U^\top$ and
effective resistance $R_e^{\rm eff}=\mathbf b_e^\top L^\dagger\mathbf b_e$.
The quadratic output inequality gives
\begin{equation}
 |\beta_e(t_k)-\beta_e^*|^2
 \le V_k[\bar P^{-1}]_{11}R_e^{\rm eff}.
 \label{detail:eq:btd_edge_bound}
\end{equation}
Under ZOH the phase, hence each fluid occupancy, is affine on an interval.
For $0\le s\le1$ its deviation is a convex combination of endpoint
deviations, so its absolute value cannot exceed the larger endpoint
absolute value. This covers all physical time, including the initial
interval from the zero state to $x_1$. It would not apply unchanged to
arbitrary delayed integer arrival events.

\begin{btdcorollary}[Buffer safety within the reference model]
\label{detail:cor:btd_buffer_safety}
Suppose the compatible centered initialization and the invariant-domain
conditions above hold with $V_1\le R$. For each receiving edge define
\begin{align}
 m_e&=\min\{\beta_e^*,C-\beta_e^*\},\nonumber\\
 B_e(R)&=\sqrt{R[\bar P^{-1}]_{11}R_e^{\rm eff}}.
 \label{detail:eq:btd_buffer_headroom_condition}
\end{align}
If $0<\beta_e^*<C$ and $B_e(R)<m_e$ on every receiving edge,
then the unconstrained fluid execution satisfies
$0<\beta_e(t)<C$ for all $t\ge0$. In addition,
$\beta_e(t)-\beta_e^*\to0$ as $t\to\infty$.
\end{btdcorollary}
\begin{proof}
At $t_k$, the squared deviation is bounded by
$V_k[\bar P^{-1}]_{11}R_e^{\rm eff}$ and hence by $B_e(R)^2$.
The initial deviation is zero. On $[t_k,t_{k+1}]$, the deviation is a
convex combination of its endpoint deviations, so its absolute value
is at most $B_e(R)$. The strict headroom inequalities give the lower
and upper capacity inequalities. Finally $V_k\le\rho^{k-1}V_1\to0$
and the same interpolation argument imply convergence at every
physical time, not only at samples.
\end{proof}

Here $B_e(R)<m_e$ is a sufficient condition. If it fails, the certificate
is inconclusive; failure of this conservative test does not prove a
buffer overflow or underflow. The condition also concerns the
unconstrained flow. It does not identify a protection-induced stall
with successful control convergence.

\subsubsection{Frequency-difference safety and convergence}
Define the physical frequency spread by
\begin{equation}
 \Delta_\omega(t)=\max_i\omega_i(t)-\min_i\omega_i(t).
 \label{detail:eq:btd_frequency_spread}
\end{equation}
A small frequency spread means close clock frequencies; it does not
require the phases or their integer counts to coincide. For any pair
$i\ne j$, not necessarily adjacent, define
\begin{equation}
 R_{ij}^{\rm eff}=(\mathbf e_i-\mathbf e_j)^\top L^\dagger
                     (\mathbf e_i-\mathbf e_j),
 \qquad R_{\max}^{\rm eff}=\max_{i\ne j}R_{ij}^{\rm eff},
 \label{detail:eq:btd_pair_resistance}
\end{equation}
where $\mathbf e_i$ is the $i$th unit vector. Connectedness makes these
quantities finite. During $[t_k,t_{k+1})$, $\omega_i(t)=
\omega_{\rm ref}+w_{i,k}$, while $v_{k+1}=Hw_k$. Therefore
\begin{align}
 |\omega_i(t)-\omega_j(t)|^2
 &=H^{-2}|(\mathbf e_i-\mathbf e_j)^\top v_{k+1}|^2\nonumber\\
 &\le H^{-2}V_{k+1}[\bar P^{-1}]_{22}R_{ij}^{\rm eff}\nonumber\\
 &\le H^{-2}R\rho^k[\bar P^{-1}]_{22}R_{ij}^{\rm eff}.
 \label{detail:eq:btd_pair_frequency_bound}
\end{align}
The factor $V_{k+1}$, rather than $V_k$, is essential: $v_k$ represents
the frequency held during the previous interval. The case $k=0$ is
covered by the explicitly constructed bootstrap state $y_1$.
At switching times the same bound applies to the newly held frequency;
the previous interval bounds the left-hand limit.

\begin{btdcorollary}[A prescribed frequency-spread threshold]
\label{detail:cor:btd_frequency_safety}
Under the invariant-domain conditions, let $\Delta_{\rm safe}>0$ be a
specified engineering tolerance on pairwise frequency differences. If
\begin{equation}
 D(R):=\frac1H\sqrt{R[\bar P^{-1}]_{22}R_{\max}^{\rm eff}}
       <\Delta_{\rm safe},
 \label{detail:eq:btd_frequency_threshold_condition}
\end{equation}
then for every $t\in[t_k,t_{k+1})$, $k\ge0$,
\begin{equation}
 \Delta_\omega(t)\le D(R)\rho^{k/2}<\Delta_{\rm safe},
 \qquad \Delta_\omega(t)\longrightarrow0.
 \label{detail:eq:btd_frequency_decay}
\end{equation}
\end{btdcorollary}
\begin{proof}
The spread equals $\max_{i,j}|\omega_i-\omega_j|$. Maximizing
\eqref{detail:eq:btd_pair_frequency_bound} and taking square roots proves
the bound. Since $0<\rho<1$, the bound tends to zero.
\end{proof}

The threshold is a user-specified performance or hardware tolerance;
logical synchrony does not by itself prescribe its numerical value.
This proof bounds the transient as well as the limiting disagreement.
It neither proves a smaller Fuzzy frequency peak than PI nor adds a
guarantee for arbitrary delayed or stochastic executions. Frequency
spread also cannot certify positive absolute frequency: a collection of
equal nonpositive rates would have zero spread. The common frequency
must therefore be bounded separately.

\subsubsection{The unprojected common frequency}
Let $\bar w_k=n^{-1}\mathbf1^\top w_k$ be the common-frequency offset.
Put $\alpha_c=(\overline\alpha+\underline\alpha)/2$ and
$s_\alpha=(\overline\alpha-\underline\alpha)/2$. With
$g=(\kappa_I,k_p)^\top$, averaging~\eqref{detail:eq:btd_d_command_increment}
and using $\mathbf1^\top U=0$ gives
\begin{equation}
 \bar w_k-\bar w_{k-1}
 =-\frac1n\mathbf1^\top(D_{\alpha,k}-\alpha_c I_n)
 U\Lambda^{1/2}(\kappa_I\xi_k+k_p\eta_k).
 \label{detail:eq:btd_d_common_increment}
\end{equation}
The second component of $g$ is $k_p$, not $Hk_p$, since $\eta$ already
contains $H$. Bounding each factor yields
\begin{equation}
 |\bar w_k-\bar w_{k-1}|
 \le s_\alpha\sqrt{\frac{\mu_n}{n}g^\top\bar P^{-1}g}\sqrt{V_k}.
 \label{detail:eq:btd_d_common_step_bound}
\end{equation}
The geometric series in $\sqrt{\rho}$ is summable. Hence
\begin{align}
 \sup_k|\bar w_k|
 &\le |\bar w_0|+
 \frac{s_\alpha\sqrt{(\mu_n/n)R g^\top\bar P^{-1}g}}{1-\sqrt\rho}
 \nonumber\\
 &\le |\bar w_0|+\frac{2s_\alpha}{1-\rho}
 \sqrt{\frac{\mu_n}{n}R g^\top\bar P^{-1}g}=W_c.
 \label{detail:eq:btd_clock_bounds}
\end{align}
The last inequality follows from
$(1-\sqrt\rho)^{-1}=(1+\sqrt\rho)/(1-\rho)\le2/(1-\rho)$.
Absolute summability also proves convergence of $\bar w_k$. Combined with
vanishing disagreement, all node frequencies converge to one value
$\omega_\infty=\omega_{\rm ref}+\lim_k\bar w_k$.
For Fuzzy PI that limit need not be the mean initial uncontrolled frequency.

\subsubsection{Closing actuator saturation and positive frequency}
The $i$th disagreement-frequency output satisfies
\begin{equation}
 |w_{i,k}-\bar w_k|\le W_d,
 \qquad W_d=\max_i\frac1H
 \sqrt{R[\bar P^{-1}]_{22}(L^\dagger)_{ii}}.
 \label{detail:eq:btd_d_node_frequency}
\end{equation}
Applying the state bound to $y_{k+1}$ covers $w_k$; $y_1$ covers $w_0$.
As $u=w-\delta\omega$, we obtain
\begin{equation}
 |u_{i,k}|\le W_c+W_d+\|\delta\omega\|_\infty,
 \qquad \omega_i(t)\ge\omega_{\rm ref}-W_c-W_d.
 \label{detail:eq:btd_actuator_bounds}
\end{equation}
First derive these bounds for the unconstrained trajectory. If the first
is strictly below $u_{\max}$ and the second is positive, the limiter is
never reached, so the constrained execution coincides with that trajectory.
This continuation argument closes the assumption of no clipping. The
buffer bound must likewise fit within physical capacity; protecting a
buffer by blocking execution is a different outcome.

Together, the following sufficient conditions connect the controller
and network parameters to the required properties:
\begin{align}
 &V_1\le R,\qquad R\Psi_i\le r^2\quad\text{for every }i,\nonumber\\
 &0<\beta_e^*<C\quad\text{for every }e,\nonumber\\
 &R[\bar P^{-1}]_{11}R_e^{\rm eff}<m_e^2
       \quad\text{for every }e,\nonumber\\
 &R[\bar P^{-1}]_{22}R_{\max}^{\rm eff}
       <H^2\Delta_{\rm safe}^2,\nonumber\\
 &W_c+W_d+\|\delta\omega\|_\infty<u_{\max},\nonumber\\
 &\omega_{\rm ref}-W_c-W_d>0,
 \label{detail:eq:btd_joint_sufficient_conditions}
\end{align}
provided the endpoint inequalities of Lemma~\ref{detail:lem:btd_endpoint} hold.
The frequency-spread inequality is required only when a specific threshold
is an objective; convergence alone does not require choosing that threshold.
All quantities are evaluated for one consistent model, initialization
and parameter set. In particular, changing $H$, $h$, $C$, node degree,
controller gains or topology can change more than one inequality.
These are sufficient certification tests, not a characterization of the
largest safe set or necessary stability boundaries.

\subsection{The prescribed triangle certificate}

Take the three-node complete graph, giving
$\mu_2=\mu_3=3$, $R_e^{\rm eff}=2/3$ and
$(L^\dagger)_{ii}=2/9$. The parameters are
\begin{align}
 H=h&=\frac{30000000}{29999},&k_p&=\frac1{500000},&k_i&=\frac1{20000000000},\nonumber\\
 K_e&=\frac1{10},&K_{\Delta e}&=\frac{29999}{7500},&K_u&=-\frac{768}{18749375},\nonumber\\
 C&=2048,&\beta^*_e&=\beta_e(0)=1024,&\lambda_e&=1024.
 \label{detail:eq:btd_example_gains}
\end{align}
Frequencies in this numerical example are normalized, and physical time is measured in reference ticks. Use the uncontrolled frequency vector
\[
 \omega^u=(1.0002,0.9998,0.9999)^\top
\]
and its average as
$\omega_{\rm ref}=29999/30000$, so
$\delta\omega=(7,-5,-2)^\top/30000$ and $\bar w_0=0$.
At the compatible startup~\eqref{detail:eq:btd_bootstrap}, use
\begin{equation}
 \bar P=\frac1{50000}\begin{bmatrix}9&181\\181&49629\end{bmatrix},
 \quad\rho=\frac{1599}{1600},\quad r=1.
 \label{detail:eq:btd_witness}
\end{equation}
Here $\det\bar P=4139/25000000>0$ and the leading diagonal entry is
positive. Since the sum of its four entries is one,
\begin{equation}
 R_0=V_1=H^2\delta\omega^\top L\delta\omega
     =\frac{234000000}{899940001}.
 \label{detail:eq:btd_d_R0}
\end{equation}
This exact fraction is used in all inequalities. The displayed approximation
$0.2600173342$ is slightly below the exact value and is not substituted as
an upper energy radius.

The sector is $[5/6,2]$, hence the spectral endpoints are $5/2$ and 6.
For $D_\nu=\rho\bar P-B(\nu)^\top\bar P B(\nu)$, exact rational
arithmetic verifies positive leading minor and determinant at both ends:
\begin{center}\small
\begin{tabular}{@{}lll@{}}\toprule
$\nu$ & $[D_\nu]_{11}$ (approx.) & $\det D_\nu$ (approx.)\\\midrule
$5/2$ & $8.2193625299\times10^{-7}$ & $7.8876135318\times10^{-11}$\\
$6$ & $2.0776400066\times10^{-6}$ & $7.3380136479\times10^{-9}$\\
\bottomrule\end{tabular}\end{center}
The accompanying machine-readable certificate stores the exact fractions;
the decimals are only a readable display, not rounded tests of positivity.
Using $R_0$ gives the following conservative, outward bounds:
\begin{center}\small
\begin{tabular}{@{}ll@{}}\toprule
Quantity & Certified bound\\\midrule
$\sup_{i,k}(|X_{i,k}|+|Y_{i,k}|)$ & $<0.171455<1$\\
Coarse edge excursion & $<32.237442$ frames\\
$\sup_k|\bar w_k|$ & $<0.003682181$\\
$\sup_{i,k}|w_{i,k}-\bar w_k|$ & $<0.000250634$\\
$\sup_{i,k}|u_{i,k}|$ & $<0.004166148<0.01$\\
$\inf_{i,t}\omega_i(t)$ & $>0.99603385$\\
\bottomrule\end{tabular}\end{center}
The coarse excursion already fits within the 1024-frame lower and upper
headroom. The frequency bounds use the exact rational-square-root
enclosures, not a sum of prematurely rounded display values.

\subsection{The exact PI all-time baseline}

For $J\in\{\mathrm{PI},\mathrm F\}$ define
\begin{equation}
 B_J=\sup_{t\ge0}\max_e|\beta_e^J(t)-1024|.
 \label{detail:eq:btd_excursion}
\end{equation}
The supremum varies time and receiving edges, while the initialization
and uncontrolled-frequency vector remain fixed. For the same plant, initialization and gains, the PI baseline has
$D_\alpha=I_n$. Averaging its increments shows that the mean frequency is
preserved. On the triangle, both disagreement modes have eigenvalue 3.
Define $\chi_0=0$, $\chi_1=1$ and
\begin{equation}
 \chi_{k+1}=(2-3a-3b)\chi_k-(1-3a)\chi_{k-1}.
 \label{detail:eq:btd_pi_scalar}
\end{equation}
To derive it, use $v_k=x_k-x_{k-1}$ in the PI state recurrence:
\[
 x_{k+1}=(2I_n-(a+b)L)x_k-(I_n-aL)x_{k-1}.
\]
Since $L\delta\omega=3\delta\omega$ and $x_1=H\delta\omega$,
induction gives $x_k=H\chi_k\delta\omega$. With
$\Delta_0=\max_i\omega_i^u-\min_i\omega_i^u=1/2500$,
\begin{equation}
 \max_e|\beta_e(t_k)-1024|=H\Delta_0|\chi_k|.
 \label{detail:eq:btd_d_pi_deviation}
\end{equation}
This exact rational sequence is enumerated through $N=200$. Its unique
prefix maximum occurs at $k=111$, with
\begin{equation}
 23.40370130812701\le
 \max_{0\le k\le N}H\Delta_0|\chi_k|\le23.40370130812702.
 \label{detail:eq:btd_d_pi_peak}
\end{equation}
It is not yet an all-time assertion until the tail is bounded.

At $N=200$ put $\delta\chi_N=\chi_N-\chi_{N-1}$, the scalar response difference. The exact state energy is
\begin{equation}
 V_N=R_0\{\bar P_{11}\chi_N^2+2\bar P_{12}\chi_N\delta\chi_N+\bar P_{22}(\delta\chi_N)^2\}.
 \label{detail:eq:btd_d_pi_VN}
\end{equation}
The common certificate also covers PI's scalar value $\nu=3$. Thus for
all $k\ge N$, the remaining buffer excursion is at most
$\sqrt{(2/3)[\bar P^{-1}]_{11}V_N}$, with an outward upper bound
less than 17.958057 frames. This is strictly below the attained prefix
maximum. ZOH interpolation therefore makes~\eqref{detail:eq:btd_d_pi_peak}
the exact all-physical-time $B_{\rm PI}$ enclosure, reached at time $111H$.

\subsection{The Fuzzy interval prefix and contracting tail}

A direct floating-point trajectory can locate a candidate peak but does
not enclose the exact one. The certificate instead propagates intervals
for all components of $x_k$ and $v_k$ on a grid $S^{-1}=10^{-60}$.
A rational number $q_{\rm rat}$ is enclosed by
\[
 \left[\frac{\lfloor Sq_{\rm rat}\rfloor}{S},\frac{\lceil Sq_{\rm rat}\rceil}{S}\right].
\]
Each multiplication uses the extrema of the four endpoint products,
rounded outward; division uses a strictly positive enclosed denominator
and outward-rounded endpoint quotients. Addition and subtraction use
endpoint arithmetic. These operations use integers and rational comparisons,
not floating-point rounding assumptions.

For the central formula, $m=\min(|X|,|Y|)$ is enclosed by interval
absolute values and endpoint minima. Its interaction term can be written
$J=\min(X,Y)$ for two positive inputs, $J=\max(X,Y)$ for two negative
inputs and zero otherwise. It is coordinatewise nondecreasing, including
across the axes; its lower and upper corner values enclose it on any
input rectangle. The positive denominator $1+2m/5$ then permits a valid
interval evaluation of the exact map. The invariant-domain proof ensures
that this central formula is applicable to the exact trajectory.

Propagating the genuine bootstrap state through state $N=200$ gives a
prefix maximum enclosed by
\begin{equation}
 [17.314056817962,\;17.314056817963]\quad\text{frames}.
 \label{detail:eq:btd_d_fuzzy_prefix}
\end{equation}
The prefix retains all node coordinates, so common drift is not silently
discarded. On the triangle the quadratic energy is equivalently
\begin{equation}
 \begin{aligned}
 V_k=\sum_{i<j}\Bigl\{&\bar P_{11}(x_{i,k}-x_{j,k})^2\\
 &+2\bar P_{12}(x_{i,k}-x_{j,k})(v_{i,k}-v_{j,k})\\
 &+\bar P_{22}(v_{i,k}-v_{j,k})^2\Bigr\}.
 \end{aligned}
 \label{detail:eq:btd_d_pair_energy}
\end{equation}
Evaluating this expression outward gives $\overline V_N\ge V_N$.
For all later samples,
\begin{equation}
 |\beta_e(t_k)-1024|^2
 \le\frac23[\bar P^{-1}]_{11}\rho^{k-N}\overline V_N.
 \label{detail:eq:btd_d_fuzzy_tail}
\end{equation}
For either controller, an equivalent tail bound that makes the
finite-prefix/infinite-tail decomposition explicit is
\begin{equation}
 T_N=\sqrt{\frac23[\bar P^{-1}]_{11}\overline V_N}.
 \label{detail:eq:btd_tail}
\end{equation}
For PI one can take $\overline V_N=V_N$ exactly; for Fuzzy,
$\overline V_N$ is an outward interval upper bound. Once $T_N$ is below
the certified prefix lower bound, no later sample can exceed the prefix
maximum. This closes the infinite time horizon without selecting an
arbitrarily long simulation duration.

An outward upper bound for this whole tail is
$12.29271005372874$ frames, strictly below the prefix lower bound.
Including $x_0=0$ and intersample interpolation gives
$B_{\rm F}<17.315$ on the entire physical time axis. The exact Fuzzy
peak is more tightly enclosed by~\eqref{detail:eq:btd_d_fuzzy_prefix}; the
rounded 17.315 value is used for a simple, conservative comparison.

The enclosures used in the comparison are collected in
Table~\ref{detail:tab:btd_prefix_tail}. The prefix intervals are outward
enclosures of actual peak values; the tail entries are upper bounds.
Their different roles are essential for the strict comparison.
\begin{table}[t]
\caption{Certified excursions in the prescribed fluid reference (frames).}
\label{detail:tab:btd_prefix_tail}
\centering\small
\begin{tabular}{@{}lll@{}}
\toprule
Controller & Prefix peak enclosure & Tail upper bound\\
\midrule
PI & $[23.40370130812701,$ & $17.95805698665804$\\
   & $\phantom{[}23.40370130812702]$ & \\
Fuzzy & $[17.314056817962,$ & $12.29271005372874$\\
      & $\phantom{[}17.314056817963]$ & \\
\bottomrule
\end{tabular}
\end{table}

\subsection{Strict comparison, engineering meaning and scope}
\label{detail:subsec:btd_comparative}

\begin{btdtheorem}[Strict headroom benefit for a prescribed initialization]
\label{detail:thm:btd_headroom}
For the triangle, parameters and compatible initialization above,
initialize the PI sample sum, the Fuzzy accumulator and the stored
previous percentage error to zero. Both reference executions keep
positive physical frequencies, remain strictly inside the buffer
interval $(0,2048)$ and never activate the $0.01$ correction limiter.
Buffer deviations and all pairwise frequency differences converge to
zero. Their all-time excursions satisfy
\begin{equation}
 B_{\rm F}<17.315<23.40370130812701\le B_{\rm PI}.
 \label{detail:eq:btd_strict_comparison}
\end{equation}
For the minimum lower/upper headroom $M_J=1024-B_J$,
\begin{equation}
 M_{\rm F}>1006.685,\qquad
 M_{\rm F}-M_{\rm PI}>6.08870130812701\quad\text{frames}.
 \label{detail:eq:btd_margin_improvement}
\end{equation}
\end{btdtheorem}
\begin{proof}
The endpoint inequalities and the radius test certify the invariant
nonlinear input region. The buffer, correction and absolute-frequency
bounds establish capacity safety, absence of clipping and positive
frequencies. The invariant disagreement energy contracts, so the buffer
and frequency-difference conclusions follow. PI is included by taking
$\alpha_i=1$, which belongs to the same sector; PI also preserves its
initial mean frequency. The certified finite-prefix enclosures and the
tail bounds imply that both all-time excursions are attained in the
prefix. The Fuzzy prefix upper bound is below $17.315$, whereas PI
attains the enclosed peak at sample 111. Intersample affine occupancy
rules out a missed larger fluid peak. Finally, each edge has a reverse
edge with the negative occupancy deviation, so the worst lower/upper
headroom is exactly $1024-B_J$. Subtracting the strict excursion bounds
gives the stated gain in headroom.
\end{proof}

The comparison uses a certified Fuzzy upper bound and an attained PI
peak. Merely comparing two conservative upper bounds would not establish
the strict ordering of the actual excursions.

Both reserves are already large. The comparison establishes a conditional
resource-margin benefit, not a demonstrated storage bottleneck or hardware
cost reduction at this operating point.

The supremum is over time and edges for one prescribed uncontrolled-frequency
vector. It is not an expanded safe set for all 400-ppm patterns, a guarantee
against arbitrary latency perturbations or superiority over independently
tuned PI. For a different capacity, $\gamma_i=100/(d_iC)$ changes; keeping
the Fuzzy gains unchanged generally changes the frame-count nonlinear
map. A new capacity certificate requires a fresh derivation or a verified
map-preserving normalization change. Physical occupancy constraints also
do not replace the delayed runtime's in-flight and credit/readiness rules.

This is a computer-assisted mathematical certificate with reproducible
integer/rational checks. It is not proof-assistant verification of the
Julia implementation. The actual asynchronous, delayed, quantized cases
belong in the experimental section. This division preserves the distinction
between a theorem about the declared reference and empirical evidence
about Callisto.

